\documentclass{article}
\usepackage[utf8]{inputenc}
\usepackage{geometry}
\usepackage{pdfpages}
\usepackage{mathtools}
\usepackage{array}
\usepackage{filecontents}
\usepackage{nicematrix}
\usepackage{bodeplot}
\graphicspath{{images/}}
\geometry{
	a4paper,
	total={170mm,257mm},
	left=20mm,
	top=20mm,
}
\usepackage{graphicx}
\usepackage{titling}
\usepackage[american,RPvoltages]{circuiTikz}
\usepackage{tikz, tikz-cd}
\usepackage{pgfplots}
\usepackage{siunitx}
\usepackage{float}
\usepackage{enumitem}
\usepackage{amsmath,amsfonts,stmaryrd,amssymb} % Math packages
\usepackage{schemabloc}
\usepackage{tcolorbox}
\usepackage{xcolor}
\usetikzlibrary{positioning}
\usepackage[scr]{rsfso}
\newcommand*{\Scale}[2][4]{\scalebox{#1}{$#2$}}%
\newcommand*{\Resize}[2]{\resizebox{#1}{!}{$#2$}}%

\newcommand{\Laplace}{\mathscr{L}}

\title{Análisis de red divisora de voltaje con potenciómetro}
\author{Rodrigo}
\date{Agosto 2024}

\usepackage{fancyhdr}
\fancypagestyle{plain}{%  the preset of fancyhdr 
	\fancyhf{} % clear all header and footer fields
	\fancyfoot[L]{\thedate}
	\fancyhead[L]{Análisis de circuito con OpA}
	\fancyhead[R]{\theauthor}
}
\makeatletter
\def\@maketitle{%
	\newpage
	\null
	\vskip 1em%
	\begin{center}%
		\let \footnote \thanks
		{\LARGE \@title \par}%
		\vskip 1em%
		%{\large \@date}%
	\end{center}%
	\par
	\vskip 1em}
\makeatother

\usepackage{lipsum}  
\usepackage{cmbright}
\renewcommand{\figurename}{Figura}
%---------------------------------------------------
\tikzset{flipflop myRS/.style={flipflop,
		flipflop def={t1=R, t3=S, t6=Q, t4={\ctikztextnot{Q}}, tu={\ctikztextnot{RST}}, nu=1, n4=1 }}
}

\tikzset{flipflop myFFRS/.style={flipflop,
		flipflop def={t1=R, t3=S, t6=Q, t4={\ctikztextnot{Q}}, n4=1 }}
}
%---------------------------------------------------

\makeatletter
\newcommand*{\underarrow}{\def\@underarrow{\relax}\@ifstar{\@@underarrow}{\def\@underarrow{\hidewidth}\@@underarrow}}
\newcommand*{\@@underarrow}[2][]{\underset{\@underarrow\substack{\uparrow\if\relax\detokenize{#1}\relax\else\\#1\fi}\@underarrow}{#2}}

\newcommand*{\overarrow}{\def\@overarrow{\relax}\@ifstar{\@@overarrow}{\def\@overarrow{\hidewidth}\@@overarrow}}
\newcommand*{\@@overarrow}[2][]{\overset{\@overarrow\substack{\if\relax\detokenize{#1}\relax\else#1\\\fi\downarrow}\@overarrow}{#2}}
\makeatother

\makeatletter
\newcounter{elimination@steps}
\newcolumntype{R}[1]{>{\raggedleft\arraybackslash$}p{#1}<{$}}
\def\elimination@num@rights{}
\def\elimination@num@variables{}
\def\elimination@col@width{}
\newenvironment{elimination}[4][0]
{
	\setcounter{elimination@steps}{0}
	\def\elimination@num@rights{#1}
	\def\elimination@num@variables{#2}
	\def\elimination@col@width{#3}
	\renewcommand{\arraystretch}{#4}
	\start@align\@ne\st@rredtrue\m@ne
}
{
	\endalign
	\ignorespacesafterend
}
\newcommand{\eliminationstep}[2]
{
	\ifnum\value{elimination@steps}>0\sim\quad\fi
	\left[
	\ifnum\elimination@num@rights>0
	\begin{array}
		{@{}*{\elimination@num@variables}{R{\elimination@col@width}}
			|@{}*{\elimination@num@rights}{R{\elimination@col@width}}}
		\else
		\begin{array}
			{@{}*{\elimination@num@variables}{R{\elimination@col@width}}}
			\fi
			#1
		\end{array}
		\right]
		& 
		\begin{array}{l}
			#2
		\end{array}
		&%                                    moved second & here
		\addtocounter{elimination@steps}{1}
	}
	\makeatother
	
	\makeatletter
	\renewcommand*\env@matrix[1][*\c@MaxMatrixCols c]{%
		\hskip -\arraycolsep
		\let\@ifnextchar\new@ifnextchar
		\array{#1}}
	\makeatother
	
\setlength\parindent{0pt}
\begin{document}
	
	\maketitle
	
	\noindent\begin{tabular}{@{}ll}
		Author & \theauthor\\

	\end{tabular}

	
	\section*{Análisis}
Sea la red divisora de voltaje simétrica, formada por dos resistencias de igual valor $R$ y un potenciómetro $R_p$, determine el valor de $R$ para cualquier valor especificado de $V_o$

\begin{figure}[H]
	\centering	
	\begin{tikzpicture}[scale=0.8]
		%\ctikzset{resistor = european}
		\ctikzset{bipoles/thickness=4}
		\ctikzset{monopoles/vcc/arrow={Stealth[width=6pt,length=9pt]}}
		\ctikzset{monopoles/vee/arrow={Stealth[width=6pt,length=9pt]}}
		\draw (0,0) node[vcc](VCC){$V_{CC}$} to[R, l_=$R$] ++(0,-3) node[potentiometershape, anchor=left, rotate=-90](P1){};
		\draw (P1.right) to[R, l_=$R$] ++(0,-3) node[vee](VEE){$V_{EE}$};
		\draw (P1.wiper) node[xshift=-11mm]{$R_p$} to[short, -o] ++(2,0) node[right]{$V_o$};
	\end{tikzpicture}	
	\caption{Red divisora de voltaje}	
	\label{fig:figura1}
\end{figure}

Redibujamos el circuito de la figura 1 como se muestra a continuación, reemplazando el potenciómetro por su modelo equivalente y colocando las fuentes de tensión apropiadas

\begin{figure}[H]
	\centering	
	\begin{tikzpicture}[scale=0.8]
		%\ctikzset{resistor = european}
		\ctikzset{bipoles/thickness=4}
		\ctikzset{monopoles/vcc/arrow={Stealth[width=6pt,length=9pt]}}
		\ctikzset{monopoles/vee/arrow={Stealth[width=6pt,length=9pt]}}
		\draw (0,0) coordinate (o1) to[R, l_=$R$] ++(0,-3) to[R, l_=$(1 - \alpha)R_p$] ++(0,-3) coordinate (p1) to[R, l_=$\alpha R_p$] ++(0,-3) to[R, l_=$R$] ++(0,-3) coordinate (o2);
		\draw (p1) to[short, *-o] ++(2,0) node[right]{$V_o$} ; 
		\draw (o1) to[short] ++(-6,0) coordinate (p2);
		\draw (p2) to[V, invert, l_=$V_{CC}$] (p2 |- p1);
		\draw (o2) to[short] ++(-6,0) coordinate (p3);
		\draw (p3) to[V, l=$V_{EE}$] (p3 |- p1) coordinate (ground);
		\draw (ground) to[short, *-] ++(2,0) node[ground]{};
	\end{tikzpicture}	
	\caption{Red divisora de voltaje}	
	\label{fig:figura2}
\end{figure}

Aplicamos el principio de superposición para encontrar el voltaje respecto a cada fuente. Sea el voltaje $V_{o_1}$ debido a $V_{CC}$ y con $V_{EE}$ desconectada.

\begin{figure}[H]
	\centering	
	\begin{tikzpicture}[scale=0.8]
		%\ctikzset{resistor = european}
		\ctikzset{bipoles/thickness=4}
		\ctikzset{monopoles/vcc/arrow={Stealth[width=6pt,length=9pt]}}
		\ctikzset{monopoles/vee/arrow={Stealth[width=6pt,length=9pt]}}
		\draw (0,0) to[V] ++(0,6) to[R, l=$R$] ++(3,0) to[R, l=$(1 - \alpha)R_p$] ++(3,0) coordinate (p1) to[R, l=$\alpha R_p$] ++(0,-3) to[R, l=$R$] ++(0,-3) coordinate (m1) to[short] (0,0);
		\draw (p1) to[short, *-o] ++(2,0) node[right]{$V_{o_1}$};
		\coordinate (M) at ($(0,0)!0.5!(m1)$);
		\draw (M) to[short, *-] ++(0,-0.1) node[ground]{};
	\end{tikzpicture}	
	\caption{Red divisora de voltaje}	
	\label{fig:figura3}
\end{figure}

Mediante divisor de tensión obtenemos

\begin{align}
	V_{o_1} = \frac{R + \alpha R_p}{R + (1 - \alpha)R_p + \alpha R_p + R} V_{CC} = \frac{R + \alpha R_p}{2R + R_p}V_{CC}
\end{align}

Ahora encontramos el voltaje $V_{o2}$ debido a $V_{EE}$, desconectando ahora la fuente de tensión $V_{CC}$

\begin{figure}[H]
	\centering	
	\begin{tikzpicture}[scale=0.8]
		%\ctikzset{resistor = european}
		\ctikzset{bipoles/thickness=4}
		\ctikzset{monopoles/vcc/arrow={Stealth[width=6pt,length=9pt]}}
		\ctikzset{monopoles/vee/arrow={Stealth[width=6pt,length=9pt]}}
		\draw (0,0) to[V, invert] ++(0,6) to[R, l=$R$] ++(3,0) to[R, l=$\alpha R_p$] ++(3,0) coordinate (p1) to[R, l=$(1 - \alpha) R_p$] ++(0,-3) to[R, l=$R$] ++(0,-3) coordinate (m1) to[short] (0,0);
		\draw (p1) to[short, *-o] ++(2,0) node[right]{$V_{o_2}$};
		\coordinate (M) at ($(0,0)!0.5!(m1)$);
		\draw (M) to[short, *-] ++(0,-0.1) node[ground]{};
	\end{tikzpicture}	
	\caption{Red divisora de voltaje}	
	\label{fig:figura4}
\end{figure}

\begin{align}
	V_{o_2} = -\frac{R + (1 - \alpha) R_p}{R + (1 - \alpha)R_p + \alpha R_p + R} V_{CC} = -  \frac{R + (1 - \alpha) R_p}{2R + R_p}V_{EE}
\end{align}

Al combinar las ecuaciones (1) y (2) obtenemos el voltaje de salida completo

\begin{align}
	V_o = V_{o_1} + V_{o_2} = \frac{R + \alpha R_p}{2R + R_p}V_{CC} - \frac{R + (1 - \alpha) R_p}{2R + R_p}V_{EE}
\end{align}

sumiendo voltajes simétricos, podemos hacer que $V_{CC} = V_{EE}$, es decir

\begin{align}
	V_o = \frac{R + \alpha R_p - R - R_p + \alpha R_p}{2R + R_p} V_{CC} = \frac{R_p (2\alpha -1 )}{R_p + 2R} V_{CC}
\end{align}

Despejamos $R$ de la ecuación (4)
\begin{align}
	\frac{V_o}{V_{CC}} = \frac{R_p (2\alpha -1 )}{R_p + 2R} \nonumber \\
	R_p V_o + 2RV_o = V_{CC}R_p(2\alpha -1) \nonumber \\
	2RV_o = V_{CC}R_p(2\alpha -1) - R_p V_o
\end{align}

Es decir

\begin{align}
	R = \frac{R_p \left[ (2\alpha - 1)V_{CC} - V_o \right]}{2V_o}
\end{align}

El voltaje de salida máximo o mínimo deseado sucederá cuando el potenciómetro se encuentre en sus extremos

\end{document}